\documentclass{article}
\usepackage{karnaugh-map}
\usepackage[utf8]{inputenc}
\usepackage[table,xcdraw]{xcolor}
\title{Práctica02}
\author{Arteaga Hernández Alejandro Iabin\\Barocio Gomez Luis Fernando\\Olea Olvera Marco Ivan}
\date{}

\begin{document}

\maketitle

\section{Ejercicios}
1.-La lógica para resolver el ejercicio 1 fué, la siguente, si A es 0 entonces siempre se hace 1, si A es 1 entonces la implicación es 1 sii 
B es 1\\\\
2.-\\EJERCICIO 2
\centering
\label{my-label}
\begin{tabular}{|c|c|c|c|}
\hline
x1 & x2 & x3 & primo \\ \hline
0  & 0  & 0  & 0     \\ \hline
0  & 0  & 1  & 0     \\ \hline
0  & 1  & 0  & 1     \\ \hline
0  & 1  & 1  & 1     \\ \hline
1  & 0  & 0  & 0     \\ \hline
1  & 0  & 1  & 1     \\ \hline
1  & 1  & 0  & 0     \\ \hline
1  & 1  & 1  & 1     \\ \hline
\end{tabular}
\\
\begin{karnaugh-map}[4][2][1][$X_2X_3$][$X_1$]
        \minterms{3,7,2,5}
        \maxterms{0,1,4,6}
        \indeterminants{}
        \implicant{3}{2}
        \implicant{5}{7}
    \end{karnaugh-map}
Mintérmino := $\overline{X_1}X_2 + X_1X_3 $\\

\centering
\textbf{EJERCICIO 3}\\
\begin{tabular}{|c|c|c|c|c|}
\hline
x1 & x2 & y1 & y2 & x \textless y \\ \hline
0  & 0  & 0  & 0  & 0             \\ \hline
0  & 0  & 0  & 1  & 1             \\ \hline
0  & 0  & 1  & 0  & 1             \\ \hline
0  & 0  & 1  & 1  & 1             \\ \hline
0  & 1  & 0  & 0  & 0             \\ \hline
0  & 1  & 0  & 1  & 0             \\ \hline
0  & 1  & 1  & 0  & 1             \\ \hline
0  & 1  & 1  & 1  & 1             \\ \hline
1  & 0  & 0  & 0  & 0             \\ \hline
1  & 0  & 0  & 1  & 0             \\ \hline
1  & 0  & 1  & 0  & 0             \\ \hline
1  & 0  & 1  & 1  & 1             \\ \hline
1  & 1  & 0  & 0  & 0             \\ \hline
1  & 1  & 0  & 1  & 0             \\ \hline
1  & 1  & 1  & 0  & 0             \\ \hline
1  & 1  & 1  & 1  & 0             \\ \hline
\end{tabular}

 \begin{karnaugh-map}
        \manualterms{0,1,1,1,0,0,1,1,0,0,0,1,0,0,0,0}
        \implicant{1}{3}
        \implicantedge{7}{6}{3}{2}
        \implicantedge{3}{3}{11}{11}
    \end{karnaugh-map}
    Mintérminos := $\overline{X_1X_2} Y_2 + \overline{X_1}Y_1+\overline{X_2}Y_1Y_2 $.
    
    

\centering
\textbf{EJERCICIO 3 parte 2}\\
\begin{tabular}{|c|c|c|c|c|}
\hline
x1 & x2 & y1 & y2 & x = y \\ \hline
0  & 0  & 0  & 0  & 1     \\ \hline
0  & 0  & 0  & 1  & 0     \\ \hline
0  & 0  & 1  & 0  & 0     \\ \hline
0  & 0  & 1  & 1  & 0     \\ \hline
0  & 1  & 0  & 0  & 0     \\ \hline
0  & 1  & 0  & 1  & 1     \\ \hline
0  & 1  & 1  & 0  & 0     \\ \hline
0  & 1  & 1  & 1  & 0     \\ \hline
1  & 0  & 0  & 0  & 0     \\ \hline
1  & 0  & 0  & 1  & 0     \\ \hline
1  & 0  & 1  & 0  & 1     \\ \hline
1  & 0  & 1  & 1  & 0     \\ \hline
1  & 1  & 0  & 0  & 0     \\ \hline
1  & 1  & 0  & 1  & 0     \\ \hline
1  & 1  & 1  & 0  & 0     \\ \hline
1  & 1  & 1  & 1  & 1     \\ \hline
\end{tabular}

\begin{karnaugh-map}
        \manualterms{1,0,0,0,0,1,0,0,0,0,1,0,0,0,0,1}
    \end{karnaugh-map}
    ¡No se pueden minimizar!\\
    Mintérminos := $\overline{X_1X_2Y_1Y_2}  + \overline{X_1}X_2\overline{Y_1}Y_2+X_1\overline{X_2}Y_1\overline{Y_2}+X_1X_2Y_1Y_2 $.
    

\centering
\textbf{EJERCICIO 4}\\
\begin{tabular}{|c|c|c|c|c|c|c|}
\hline
x1 & x2 & y1 & y2 & direccion & p1 & p2 \\ \hline
0  & 0  & 0  & 0  & 0         & 0  & 0  \\ \hline
0  & 0  & 0  & 1  & 1         & 0  & 0  \\ \hline
0  & 0  & 1  & 0  & 1         & 1  & 0  \\ \hline
0  & 0  & 1  & 1  & 1         & 1  & 1  \\ \hline
0  & 1  & 0  & 0  & 0         & 0  & 1  \\ \hline
0  & 1  & 0  & 1  & 0         & 0  & 0  \\ \hline
0  & 1  & 1  & 0  & 1         & 0  & 1  \\ \hline
0  & 1  & 1  & 1  & 1         & 1  & 0  \\ \hline
1  & 0  & 0  & 0  & 0         & 1  & 0  \\ \hline
1  & 0  & 0  & 1  & 0         & 0  & 1  \\ \hline
1  & 0  & 1  & 0  & 0         & 1  & 0  \\ \hline
1  & 0  & 1  & 1  & 1         & 0  & 1  \\ \hline
1  & 1  & 0  & 0  & 0         & 1  & 1  \\ \hline
1  & 1  & 0  & 1  & 0         & 1  & 0  \\ \hline
1  & 1  & 1  & 0  & 0         & 0  & 1  \\ \hline
1  & 1  & 1  & 1  & 0         & 0  & 1  \\ \hline
\end{tabular}
La reducción para saber si es mayor y saber si tengo que subir o bajar, ya la he hecho, por lo que solamente haré la reducción de la diferencia de pisos
\begin{karnaugh-map}
        \manualterms{0,0,1,1,0,0,0,1,1,0,0,0,1,1,0,0}
        \implicant{3}{2}
        \implicant{3}{7}
        \implicant{12}{13}
        \implicantedge{8}{8}{10}{10}
\end{karnaugh-map}
\begin{karnaugh-map}
        \manualterms{0,1,0,1,1,0,1,0,0,1,0,1,1,0,1,0}
         \implicantedge{1}{3}{9}{11}
        \implicantedge{4}{12}{6}{14}
\end{karnaugh-map}
 Mintérminos digito 1 := $\overline{X_1X_2}Y_1  + \overline{X_1}Y_1Y_2+X_1\overline{X_2Y_2}+X_1X_2\overline{Y_1} $.\\
  Mintérminos digito 2 := $\overline{X_2}Y_2  +X_2 \overline{Y_2}$.
\raggedright
  
  
    
    
\end{document}
